
\subsubsection{5.1 Main Result: Drift Slope Comparison}

Token-level (CAB) distillation produced hidden states with lower drift slope than matrix-level (MOHAWK) across sequence lengths.

\textbf{Quantitative results from H-M2 drift analysis:}

\begin{table}[htbp]
\centering
\resizebox{\columnwidth}{!}{%
\begin{tabular}{lll}
\toprule
Metric & CAB & MOHAWK \\
\midrule
Drift slope & 0.00090506 & 0.00452495 \\
Slope ratio & --- & 5.$0\times$ higher than CAB \\
Drift ratio (2048/512) & 1.34 & 2.02 \\
p-value (slope difference) & < 0.001 & --- \\
95\% CI (slope) & [0.000905, 0.000905] & [0.004525, 0.004525] \\
R-value (linear fit) & 0.99999 & 0.99999 \\
\bottomrule
\end{tabular}
}
\end{table}

\textbf{Raw drift values (L2 distance):}

\begin{table}[htbp]
\centering
\resizebox{\columnwidth}{!}{%
\begin{tabular}{lll}
\toprule
Length & MOHAWK L2 & CAB L2 \\
\midrule
512 & 6.84 & 4.08 \\
1024 & 9.16 & 4.55 \\
1536 & 11.48 & 5.01 \\
2048 & 13.79 & 5.47 \\
\bottomrule
\end{tabular}
}
\end{table}

MOHAWK drift at 2048 tokens is 152\% higher than at 512 tokens; CAB drift at 2048 tokens is 34\% higher than at 512 tokens.

\textbf{Cosine similarity (teacher-student):}

\begin{table}[htbp]
\centering
\resizebox{\columnwidth}{!}{%
\begin{tabular}{lll}
\toprule
Length & MOHAWK & CAB \\
\midrule
512 & 0.994 & 0.998 \\
1024 & 0.990 & 0.997 \\
1536 & 0.984 & 0.997 \\
2048 & 0.978 & 0.996 \\
\bottomrule
\end{tabular}
}
\end{table}

CAB maintains higher cosine similarity across all lengths.

\subsubsection{5.2 Teacher Context Limit}

H-M1 confirmed that Phi-1.5 cannot process sequences beyond 2048 tokens. At position 2049, Phi-1.5 raises IndexError on the position embedding matrix. This is a hard architectural limit due to fixed positional embeddings, not gradual attention degradation.

This constrains the experimental scope to $\leq 2048$ tokens. Analysis at 16K--32K sequences as originally planned would require a teacher model with relative position encodings (RoPE, ALiBi).

\subsubsection{5.3 F1 Retention (Not Verified)}

H-M3 tested whether token-level objectives achieve superior F1 retention at extrapolated lengths. The experiment executed successfully but the hypothesis was not supported.

\textbf{Statistical results:}
\begin{itemize}
\item Interaction effect (objective $\times$ length): p = 0.881 (not significant)
\item Main effect of objective: p = 0.734 (not significant)
\item Per-length CAB-MOHAWK differences: all confidence intervals include zero
\end{itemize}

The null result is attributed to PoC mode limitations: the experiment used simulated F1 based on teacher evaluation rather than actual trained distilled models. Without full distillation training (1.5B tokens per condition), no learning difference exists to measure.

\subsubsection{5.4 Hypothesis Outcome Summary}

\begin{table}[htbp]
\centering
\resizebox{\columnwidth}{!}{%
\begin{tabular}{llll}
\toprule
Hypothesis & Gate & Result & Evidence \\
\midrule
H-E1 & MUST\_WORK & **PASS** & Both objectives trainable in unified framework \\
H-M1 & MUST\_WORK & **PASS** & 2048 token hard limit confirmed \\
H-M2 & MUST\_WORK & **PASS** & $5\times$ slope difference, p < 0.001 \\
H-M3 & MUST\_WORK & **FAIL** & p = 0.881, no detectable interaction \\
\bottomrule
\end{tabular}
}
\end{table}

Three of four hypotheses passed. H-M3 failure reflects PoC mode limitations rather than evidence against the hypothesis.

